\scanpage{29}
\mspage{14}
% archive foliation: 28

La donnée de $(M_0,I)$ équivaut, par le §1, à la donnée d'un ens. à quatre éléments $I$, on aura $M_0\simeq V_I$, $M_0^*\simeq\mathcal P_{2,2}(I)$. \deleted{\unsure{Pour}} \inserted{D'autre} part, \unsure{\deleted{la donnée d'un}|\uline{la donnée d'un}} ens. à deux éléments $C$, et d'un torseur $S$ sous $\mathbb{F}_2^C$, équivaut à la donnée d'un carré $Q$, $S$ étant l'ens. des sommets de $Q$, $C$ l'ens de ses codiagonales.

Donc le \deleted{catégorie} \inserted{gr.ïde} du deuxième mb de (18) est équivalent à celui des triples $(I,Q;k)$, où $I$ est un ens à quatre éléments, $Q$ un carré, et $k$ un isom.
\[
(20)\qquad k:\underset{\substack{\text{ens des}\\\text{orientations}\\\text{de }I}}{\underline{\omega}(I)}
\simeq
\underset{\substack{\text{ens. des}\\\text{codiagonales}\\\text{de }Q}}{C(Q)}.
\]

On trouve donc

\uline{Corollaire 2} On a une équivalence can. de groupoïdes
\[
(21)\qquad
\underset{\mathbb{Z}/4\mathbb{Z}}{\underline{\Lambda}}\text{-modules libres de rg }2
\xrightarrow{\approx}
\text{triplets }(I,Q,k),\ \text{où}
\quad
\left\{
\begin{array}{l}
I\ \text{ens. à 4 éléments}\\
Q\ \text{carré combinatoire}\\
k:\underline{\omega}(I)\simeq C(Q).
\end{array}
\right.
\]

\uline{Corollaire 3} Si $M$ est associé à $I,Q$, on a un diagramme cartésien canonique
\[
(22)\qquad
\diagram{\begin{tikzcd}
& GL_\Lambda(M)\arrow[dl]\arrow[dr]&\\
\unsure{\mathcal{G}_I|\mathfrak S_I}\arrow[dr,"\mathrm{sg}"']&&\mathfrak D_Q\arrow[dl,"\chi_C\ (\text{opération sur l'ens des codiagonales})"]\\
&\{\pm1\}&
\end{tikzcd}}
\]
